% ============================================================
%  Chapter 4: Objectives, requirements and specification
%  Included from main.tex via % ============================================================
%  Chapter 4: Objectives, requirements and specification
%  Included from main.tex via % ============================================================
%  Chapter 4: Objectives, requirements and specification
%  Included from main.tex via % ============================================================
%  Chapter 4: Objectives, requirements and specification
%  Included from main.tex via \input{Chapters/Specification}.
%  Do NOT add \documentclass or \begin{document} here.
%  Follows Chapter 3 (Literature Review); precedes Chapter 5 (Design).
%  Scope: what the project aims to learn (objectives, RQs, hypotheses)
%  and what the system must satisfy (FR, NFR, ER, experimental
%  specification, schema, evaluation). Architecture and alternatives
%  considered belong to Chapter~\ref{chap:design}.
% ============================================================

\chapter{Objectives, requirements and specification}
\label{chap:spec}

This chapter fixes what the project aims to learn and what the
system must satisfy for that learning to count as evidence. Research
objectives and hypotheses are in \S\ref{sec:spec-objectives}, scope
in \S\ref{sec:spec-scope}, requirements in
\S\S\ref{sec:spec-fr}--\ref{sec:spec-experimental-requirements},
experimental specification in \S\ref{sec:spec-experimental}, the
schema in \S\ref{sec:spec-schema}, and evaluation criteria in
\S\ref{sec:spec-evaluation}. Every requirement traces back to a
subsection of Chapter~\ref{chap:litreview}; how each is realised
architecturally is deferred to Chapter~\ref{chap:design}. This
chapter states \textit{what} and \textit{why}; the next states
\textit{how} and \textit{what was rejected}.

% ------------------------------------------------------------
\section{Research objectives and hypotheses}
\label{sec:spec-objectives}

Three research questions follow from the gap in
\S\ref{sec:lr-gap}. Each is paired with a directional hypothesis
fixed before the run, so a null result stays informative against a
pre-registered position.

\paragraph{RQ1 --- Modality effect.} Does explicit OCR preprocessing
improve or degrade structured-extraction accuracy compared to
direct vision-mediated extraction by the same LLM, on multilingual,
out-of-distribution utility bills?

\textit{Hypothesis H1.} Vision-mediated extraction will match or
exceed OCR-mediated extraction in aggregate accuracy, with the gap
widening on non-Latin-script and low-frequency layouts, because the
OCR error-propagation ceiling identified by
Lopresti~\cite{lopresti2009ocrerrors} and Nguyen et
al.~\cite{nguyen2021postocr} binds on multilingual input and is not
shared by the vision path.

\paragraph{RQ2 --- Operational trade-off.} How do the OCR-mediated
and vision-mediated pipelines compare on per-document monetary
cost, wall-clock latency, and operational complexity?

\textit{Hypothesis H2.} Vision-mediated extraction will cost more
per document in API tokens while eliminating Tesseract language-pack
management, producing a cost-complexity trade-off rather than a
single-winner ranking, consistent with the operational-cost
observations in FastVLM~\cite{vasu2024fastvlm} and
DocVLM~\cite{nacson2024docvlm}.

\paragraph{RQ3 --- Cross-provider generality.} Does the modality
effect observed under RQ1 generalise across LLM providers from
different architectural families, or is it provider-specific?

\textit{Hypothesis H3.} The OCR-vs-vision gap will vary by provider
because structured-output support, vision-encoder resolution, and
multilingual OCR capability differ between
OpenAI~\cite{openai2024gpt4o,openai2024structured} and
Anthropic~\cite{anthropic2024claude} models. A provider-independent
gap would falsify H3.

% ------------------------------------------------------------
\section{Scope, terminology and success framing}
\label{sec:spec-scope}

\paragraph{In scope.} Template-free structured extraction from
multilingual, born-digital and scanned utility bills. Two
general-purpose LLM providers held constant across modalities.
Zero-shot prompting with temperature~$0.0$. Full-document accuracy,
cost, and latency under the null-aware evaluation of
\S\ref{sec:spec-evaluation}. Per-stage attribution of OCR-mediated
failures to OCR text or LLM reasoning.

\paragraph{Out of scope.} Specialised document architectures
(Donut~\cite{kim2022donut}, LayoutLMv3~\cite{huang2022layoutlmv3});
fine-tuning; few-shot prompting; image preprocessing beyond
rasterisation; asynchronous or batched submission; confidence
calibration; significance testing on slice accuracy (sample size
does not support it reliably). A specialised architecture, a
fine-tuning step, or a preprocessing transform would turn the
comparison into a cross-architecture or cross-training study rather
than the single-variable modality test the gap demands. These are
acknowledged as limitations in Chapter~\ref{chap:evaluation}.

\paragraph{Canonical terminology.} The literature uses several names
for each pipeline. For the rest of this report:

\begin{itemize}
\item \textbf{OCR-mediated pipeline}: document image through
Tesseract~\cite{smith2007tesseract}, then the text to the LLM in
text mode. Equivalent to ``OCR-cascade'' and ``text-mode''.
\item \textbf{Vision-mediated pipeline}: document image sent
directly to the LLM in vision mode. Equivalent to ``OCR-free'' and
``vision-mode''.
\item \textbf{Reasoning engine}: the LLM held constant across the
two pipelines within a condition.
\item \textbf{Condition}: a unique (\textit{provider},
\textit{modality}) pair; four exist in the $2\times 2$ factorial,
labelled C1--C4.
\end{itemize}

\paragraph{Success framing.} The project is not chasing a single
leaderboard number. For each RQ, success means a metric collected
under pre-specified conditions, a directional hypothesis fixed
before collection, and a mechanism-level attribution for the sign
of the result. A null finding that is attributable to a named
mechanism (OCR propagation, vision hallucination, output-format
asymmetry, provider-specific encoder limitations) is a successful
outcome; an uninterpretable headline number is not.

% ------------------------------------------------------------
\section{Functional requirements}
\label{sec:spec-fr}

Functional requirements state what the system must do. Each is
verifiable against a run artefact, an automated test, or an
evaluation metric. Priority uses a two-level MoSCoW reduction:
\textit{Must} requirements are those whose violation invalidates
the experiment; \textit{Should} requirements improve operational
usability without affecting validity.

\begin{figure}[ht]
\centering
\includegraphics[width=\linewidth]{Chapters/figures/fig-4-2-fr-table.png}
\caption[Functional requirements catalogue.]{Functional
requirements catalogue. Each requirement names both its
verification artefact and its automated test.}
\label{fig:fr}
\end{figure}

Figure~\ref{fig:fr} fixes visible behaviour, not internal structure;
how ingestion, pipelines, and evaluation are organised inside the
system is covered in Chapter~\ref{chap:design}.

% ------------------------------------------------------------
\section{Non-functional requirements}
\label{sec:spec-nfr}

Non-functional requirements state properties the system must exhibit
irrespective of any one extraction. They guard experimental validity
against engineering failure modes documented in the literature:
irreproducibility, output-mechanism asymmetry, silent configuration
drift, and runaway
cost~\cite{pineau2021reproducibility,openai2024structured}. A system
can satisfy an FR while violating an NFR; such a system produces
extractions but not an experiment.

\begin{figure}[ht]
\centering
\includegraphics[width=\linewidth]{Chapters/figures/fig-4-3-nfr-table.png}
\caption[Non-functional requirements catalogue.]{Non-functional
requirements catalogue. All eight are priority \textit{Must}.}
\label{fig:nfr}
\end{figure}

NFR1--NFR8 are all \textit{Must}; each one is a precondition for a
finding to be defensible under cold-reading scrutiny.

% ------------------------------------------------------------
\section{Experimental requirements}
\label{sec:spec-experimental-requirements}

Experimental requirements are the seven classes derived from the
rubric of \S\ref{sec:lr-controlled} and the null-handling analysis
of \S\ref{sec:lr-eval}. They differ from FRs and NFRs in being
properties of the \textit{experiment}: a system can be functionally
perfect and still fail ER1, for example by feeding the two
modalities different post-processing, and so invalidate the
comparison.

\begin{figure}[ht]
\centering
\includegraphics[width=\linewidth]{Chapters/figures/fig-4-4-er-table.png}
\caption[Experimental requirements.]{Experimental requirements,
each traced to the literature-review subsection that motivates it.}
\label{fig:er}
\end{figure}

Figure~\ref{fig:er} closes the traceability chain \textit{lit-review
subsection $\to$ ER $\to$ design decision $\to$ test $\to$ evaluation
revisit}. ER1--ER7 are all \textit{Must}; any violation invalidates
the study against the rubric this report has committed to.

% ------------------------------------------------------------
\section{Experimental specification}
\label{sec:spec-experimental}

The experiment is a $2\times 2$ factorial over input modality
(OCR-mediated, vision-mediated) and LLM provider (OpenAI,
Anthropic). The four conditions are labelled C1--C4.

\begin{figure}[ht]
\centering
\includegraphics[width=0.78\linewidth]{Chapters/figures/fig-4-1-factorial-matrix.png}
\caption[$2\times 2$ factorial matrix.]{The four conditions under
which the system is run. Rows: LLM provider. Columns: input
modality. Each cell is a complete pipeline run over the full
multilingual utility-bill dataset.}
\label{fig:factorial}
\end{figure}

Within a column, modality is fixed and the reasoning engine changes
(RQ3). Within a row, the reasoning engine is fixed and modality
changes (RQ1). All four cells are used jointly for RQ2.

Each RQ is pre-specified below: metric, baselines, hypothesis,
controls, and success criterion fixed before any run, so the
evaluation chapter only reports whether each condition met or
missed those criteria.

\paragraph{RQ1 pre-specification.} \textit{Metric:} overall,
document-level, and per-field accuracy
(\S\ref{sec:spec-evaluation}), independently for C1--C4.
\textit{Baselines:} a null-prediction baseline returning
\texttt{null} for every field, to establish the floor implied by
the both-null rule; each condition paired against its within-provider
counterpart (C1 vs.~C2, C3 vs.~C4); Shen et al.'s directional prior
that image-only MLLM input can reach parity with OCR-enhanced
setups~\cite{shen2026ocrornot}. \textit{Hypothesis:} H1.
\textit{Controls:} ER1, plus shared schema, normalisation, and
evaluation code. \textit{Success:} the within-provider modality gap
is measurable at field-level granularity and its sign is
attributable via ER5 to a named mechanism (OCR propagation, vision
hallucination, output-format reliability) rather than asserted.

\paragraph{RQ2 pre-specification.} \textit{Metric:} mean per-document
latency with the per-stage breakdown of NFR7, mean per-document
token count, and mean per-document dollar cost, each for C1--C4.
\textit{Baselines:} Tesseract-only latency as the OCR-mediated stage
ceiling; first-two-page vision-mediated cost as a controlled upper
bound on image-token input;
FastVLM~\cite{vasu2024fastvlm} throughput figures as external
reference. \textit{Hypothesis:} H2. \textit{Controls:} ER1 and
identical document set; identical run ordering where API rate limits
permit. \textit{Success:} cost, latency, and operational complexity
are reported as three separable quantities, with no single-winner
collapse without direct quantitative support.

\paragraph{RQ3 pre-specification.} \textit{Metric:} the
within-provider modality gap for each provider, and the difference
between those gaps with its sign. \textit{Baselines:} both providers
on the same dataset, prompt, normalisation, and evaluation harness,
so provider is the only uncontrolled variable across rows.
\textit{Hypothesis:} H3. \textit{Controls:} ER1, ER2, ER6.
\textit{Success:} a provider-specificity claim is accepted only
when the sign or magnitude of the modality gap matches a named
mechanism, for example native constrained
decoding~\cite{openai2024structured} versus prompt-instructed JSON,
or vision-encoder resolution for low-frequency layouts; otherwise
the null is reported.

\paragraph{Slicing.} Accuracy metrics are reported in aggregate and
sliced by \texttt{language} and \texttt{utility\_type}. Slice
accuracies are descriptive given the corpus size; formal
significance testing is deferred (out-of-scope list,
\S\ref{sec:spec-scope}).

% ------------------------------------------------------------
\section{Canonical data schema}
\label{sec:spec-schema}

The extraction target is a 12-field canonical schema
(Figure~\ref{fig:schema}). Fixing the schema is what makes FR5--FR7
objective: an unbounded key-value target has no ground truth against
which hallucination, omission, or value mismatch can be defined, and
the null-aware evaluation (\S\ref{sec:spec-evaluation}) relies on
the closed-world assumption that ``not in the schema'' is not a
valid prediction. All 12 fields are nullable; a given bill may
omit any of them.

\begin{figure}[ht]
\centering
\includegraphics[width=\linewidth]{Chapters/figures/fig-4-5-schema-table.png}
\caption[\texttt{BillExtraction} schema.]{The canonical
\texttt{BillExtraction} schema. All fields are nullable; types are
enforced by Pydantic at parse time.}
\label{fig:schema}
\end{figure}

The 12 fields are the intersection of the minimum-content
requirements mandated for electricity bills by Article~18 of the EU
electricity directive~\cite{eu2019electricity} and the stable core
fields ReceiptSense~\cite{abdallah2024receiptsense} observes across
multilingual receipts. Provider-specific fields outside that core
(per-tariff itemisation, previous-balance reconciliation) are
excluded to keep the comparison constant across issuers whose
layouts differ substantially in those areas.

\paragraph{Normalisation.} Provider-name legal-suffix stripping
(\textit{Ltd}, \textit{GmbH}, \textit{S.p.A.}); EU- and US-style
decimal-separator parsing; ISO~8601 canonicalisation of dates from
locale-specific surface forms; multilingual currency and
consumption-unit synonym mapping (Italian ``SMC'' maps to m$^3$).
Full rules are in Appendix~B.

\paragraph{Ground-truth format.} Ground truth is stored per document
as JSON:

\begin{verbatim}
{
  "document_id": "<unique key>",
  "annotated_by": "<annotator>",
  "annotation_date": "<ISO 8601>",
  "fields": {
    "provider_name": "<value or null>",
    ...
  }
}
\end{verbatim}

Fields absent from the bill are set to \texttt{null} explicitly,
distinguishing ``absent'' from ``unannotated''. The four-class null
taxonomy (\S\ref{sec:spec-evaluation}) depends on this being decided
at annotation time rather than inferred from missing keys.

% ------------------------------------------------------------
\section{Evaluation criteria}
\label{sec:spec-evaluation}

Criteria are fixed here so Chapter~\ref{chap:evaluation} reports
against a target chosen before the data is seen. Executable decision
flows for null classification and failure diagnosis are deferred to
Chapter~\ref{chap:design} \S\ref{sec:design-evalframework}.

\paragraph{Field-level accuracy.} For each field~$f$, let $C_f$,
$I_f$, $H_f$, $O_f$, $B_f$ be the counts across all documents of
correct, incorrect, hallucinated, omitted, and both-null predictions.
Field accuracy is
\(A_f = C_f / (C_f + I_f + H_f + O_f)\). The both-null count $B_f$ is
excluded from the denominator because including it would inflate
accuracy by rewarding legitimate absences, an effect flagged by
Chinchor's MUC-7 response-key vocabulary~\cite{chinchor1998muc7}
and by ExtractBench~\cite{ferguson2026extractbench}.

\paragraph{Document-level accuracy.} The fraction of documents with
zero incorrect, hallucinated, or omitted fields across all 12
targets. A single wrong field invalidates the extraction for
downstream bookkeeping, so this is the most production-relevant
measure.

\paragraph{Overall accuracy.} Field-level counts aggregated across
all 12 fields and all documents under the same both-null-excluded
denominator.

\paragraph{Four-class null taxonomy.} Every field comparison falls
into one of the five cells of Figure~\ref{fig:null-taxonomy}. The
distinction between hallucination and omission follows Dale et
al.~\cite{dale2023hallucination} and is required by ER6.

\begin{figure}[ht]
\centering
\includegraphics[width=0.85\linewidth]{Chapters/figures/fig-4-6-null-taxonomy.png}
\caption[Null-aware taxonomy.]{The null-aware evaluation taxonomy
applied to every field comparison.}
\label{fig:null-taxonomy}
\end{figure}

\paragraph{Comparison rules.} Floats (\texttt{total\_amount\_due},
\texttt{consumption\_amount}) compare with absolute tolerance
$\pm 0.01$ after normalisation. Dates compare by exact string
equality on the ISO~8601 canonical form. Strings compare on the
normalised surface form (NFC Unicode, lowercase,
whitespace-collapsed); Levenshtein similarity is computed for
diagnosis but does not contribute to accuracy.

\paragraph{Failure attribution, OCR-mediated (ER5).} For each
non-correct field: if the ground-truth value occurs as a
case-insensitive substring of the stored OCR text, the failure is
attributed to the LLM (the signal was available and it misread it);
otherwise the failure is attributed to Tesseract (the signal was
lost before the LLM saw it). The heuristic has known false-positive
and false-negative modes, especially for dates where the same digits
recur in unrelated contexts; its limitations are discussed in
Chapter~\ref{chap:design}. The heuristic supplies an automated,
reproducible, first-order attribution that the OCR-cascade
literature (\S\ref{sec:lr-ocrcascade}) does not.

\paragraph{Failure attribution, vision-mediated.} All non-correct
fields are attributed to the vision-mediated model without finer
granularity, because no intermediate artefact analogous to OCR text
exists. The asymmetry between the two conditions is itself a finding
of the study; Chapter~\ref{chap:evaluation} discusses it rather than
papering over it with a weaker heuristic.

\paragraph{Chart set (FR8).} Six canonical figures are produced by
the reporting stage: (1) overall-accuracy bar chart across
conditions; (2) per-field-accuracy heatmap; (3) per-stage timing
stacked bar chart; (4) cost-per-document bar chart;
(5) accuracy-by-language grouped bar chart; (6)
accuracy-by-\texttt{utility\_type} grouped bar chart. Each supports
one or more of RQ1--RQ3; the chart-to-RQ mapping is given in
Chapter~\ref{chap:design}.

% ------------------------------------------------------------
\section{Summary and forward bridge}
\label{sec:spec-bridge}

This chapter has fixed RQ1--RQ3 and their directional hypotheses,
scope and terminology, FR1--FR10, NFR1--NFR8, and ER1--ER7 with
upstream traceability to Chapter~\ref{chap:litreview}, the $2\times 2$
experimental matrix and per-RQ success criteria, the 12-field
schema, and the null-aware evaluation metrics and failure-attribution
protocol.

Chapter~\ref{chap:design} takes this specification as input. It
states the architectural style chosen to satisfy FR1--FR10 under
NFR1--NFR8 and ER1--ER7, names the ten design decisions that
realise those requirements, identifies the alternative considered
for each, and gives the technical reason each alternative was
rejected. The evaluation framework specified declaratively here is
given its decision-tree realisation in
Chapter~\ref{chap:design}~\S\ref{sec:design-evalframework}.
Chapter~\ref{chap:evaluation} closes the chain by revisiting every
requirement against the run evidence.
.
%  Do NOT add \documentclass or \begin{document} here.
%  Follows Chapter 3 (Literature Review); precedes Chapter 5 (Design).
%  Scope: what the project aims to learn (objectives, RQs, hypotheses)
%  and what the system must satisfy (FR, NFR, ER, experimental
%  specification, schema, evaluation). Architecture and alternatives
%  considered belong to Chapter~\ref{chap:design}.
% ============================================================

\chapter{Objectives, requirements and specification}
\label{chap:spec}

This chapter fixes what the project aims to learn and what the
system must satisfy for that learning to count as evidence. Research
objectives and hypotheses are in \S\ref{sec:spec-objectives}, scope
in \S\ref{sec:spec-scope}, requirements in
\S\S\ref{sec:spec-fr}--\ref{sec:spec-experimental-requirements},
experimental specification in \S\ref{sec:spec-experimental}, the
schema in \S\ref{sec:spec-schema}, and evaluation criteria in
\S\ref{sec:spec-evaluation}. Every requirement traces back to a
subsection of Chapter~\ref{chap:litreview}; how each is realised
architecturally is deferred to Chapter~\ref{chap:design}. This
chapter states \textit{what} and \textit{why}; the next states
\textit{how} and \textit{what was rejected}.

% ------------------------------------------------------------
\section{Research objectives and hypotheses}
\label{sec:spec-objectives}

Three research questions follow from the gap in
\S\ref{sec:lr-gap}. Each is paired with a directional hypothesis
fixed before the run, so a null result stays informative against a
pre-registered position.

\paragraph{RQ1 --- Modality effect.} Does explicit OCR preprocessing
improve or degrade structured-extraction accuracy compared to
direct vision-mediated extraction by the same LLM, on multilingual,
out-of-distribution utility bills?

\textit{Hypothesis H1.} Vision-mediated extraction will match or
exceed OCR-mediated extraction in aggregate accuracy, with the gap
widening on non-Latin-script and low-frequency layouts, because the
OCR error-propagation ceiling identified by
Lopresti~\cite{lopresti2009ocrerrors} and Nguyen et
al.~\cite{nguyen2021postocr} binds on multilingual input and is not
shared by the vision path.

\paragraph{RQ2 --- Operational trade-off.} How do the OCR-mediated
and vision-mediated pipelines compare on per-document monetary
cost, wall-clock latency, and operational complexity?

\textit{Hypothesis H2.} Vision-mediated extraction will cost more
per document in API tokens while eliminating Tesseract language-pack
management, producing a cost-complexity trade-off rather than a
single-winner ranking, consistent with the operational-cost
observations in FastVLM~\cite{vasu2024fastvlm} and
DocVLM~\cite{nacson2024docvlm}.

\paragraph{RQ3 --- Cross-provider generality.} Does the modality
effect observed under RQ1 generalise across LLM providers from
different architectural families, or is it provider-specific?

\textit{Hypothesis H3.} The OCR-vs-vision gap will vary by provider
because structured-output support, vision-encoder resolution, and
multilingual OCR capability differ between
OpenAI~\cite{openai2024gpt4o,openai2024structured} and
Anthropic~\cite{anthropic2024claude} models. A provider-independent
gap would falsify H3.

% ------------------------------------------------------------
\section{Scope, terminology and success framing}
\label{sec:spec-scope}

\paragraph{In scope.} Template-free structured extraction from
multilingual, born-digital and scanned utility bills. Two
general-purpose LLM providers held constant across modalities.
Zero-shot prompting with temperature~$0.0$. Full-document accuracy,
cost, and latency under the null-aware evaluation of
\S\ref{sec:spec-evaluation}. Per-stage attribution of OCR-mediated
failures to OCR text or LLM reasoning.

\paragraph{Out of scope.} Specialised document architectures
(Donut~\cite{kim2022donut}, LayoutLMv3~\cite{huang2022layoutlmv3});
fine-tuning; few-shot prompting; image preprocessing beyond
rasterisation; asynchronous or batched submission; confidence
calibration; significance testing on slice accuracy (sample size
does not support it reliably). A specialised architecture, a
fine-tuning step, or a preprocessing transform would turn the
comparison into a cross-architecture or cross-training study rather
than the single-variable modality test the gap demands. These are
acknowledged as limitations in Chapter~\ref{chap:evaluation}.

\paragraph{Canonical terminology.} The literature uses several names
for each pipeline. For the rest of this report:

\begin{itemize}
\item \textbf{OCR-mediated pipeline}: document image through
Tesseract~\cite{smith2007tesseract}, then the text to the LLM in
text mode. Equivalent to ``OCR-cascade'' and ``text-mode''.
\item \textbf{Vision-mediated pipeline}: document image sent
directly to the LLM in vision mode. Equivalent to ``OCR-free'' and
``vision-mode''.
\item \textbf{Reasoning engine}: the LLM held constant across the
two pipelines within a condition.
\item \textbf{Condition}: a unique (\textit{provider},
\textit{modality}) pair; four exist in the $2\times 2$ factorial,
labelled C1--C4.
\end{itemize}

\paragraph{Success framing.} The project is not chasing a single
leaderboard number. For each RQ, success means a metric collected
under pre-specified conditions, a directional hypothesis fixed
before collection, and a mechanism-level attribution for the sign
of the result. A null finding that is attributable to a named
mechanism (OCR propagation, vision hallucination, output-format
asymmetry, provider-specific encoder limitations) is a successful
outcome; an uninterpretable headline number is not.

% ------------------------------------------------------------
\section{Functional requirements}
\label{sec:spec-fr}

Functional requirements state what the system must do. Each is
verifiable against a run artefact, an automated test, or an
evaluation metric. Priority uses a two-level MoSCoW reduction:
\textit{Must} requirements are those whose violation invalidates
the experiment; \textit{Should} requirements improve operational
usability without affecting validity.

\begin{figure}[ht]
\centering
\includegraphics[width=\linewidth]{Chapters/figures/fig-4-2-fr-table.png}
\caption[Functional requirements catalogue.]{Functional
requirements catalogue. Each requirement names both its
verification artefact and its automated test.}
\label{fig:fr}
\end{figure}

Figure~\ref{fig:fr} fixes visible behaviour, not internal structure;
how ingestion, pipelines, and evaluation are organised inside the
system is covered in Chapter~\ref{chap:design}.

% ------------------------------------------------------------
\section{Non-functional requirements}
\label{sec:spec-nfr}

Non-functional requirements state properties the system must exhibit
irrespective of any one extraction. They guard experimental validity
against engineering failure modes documented in the literature:
irreproducibility, output-mechanism asymmetry, silent configuration
drift, and runaway
cost~\cite{pineau2021reproducibility,openai2024structured}. A system
can satisfy an FR while violating an NFR; such a system produces
extractions but not an experiment.

\begin{figure}[ht]
\centering
\includegraphics[width=\linewidth]{Chapters/figures/fig-4-3-nfr-table.png}
\caption[Non-functional requirements catalogue.]{Non-functional
requirements catalogue. All eight are priority \textit{Must}.}
\label{fig:nfr}
\end{figure}

NFR1--NFR8 are all \textit{Must}; each one is a precondition for a
finding to be defensible under cold-reading scrutiny.

% ------------------------------------------------------------
\section{Experimental requirements}
\label{sec:spec-experimental-requirements}

Experimental requirements are the seven classes derived from the
rubric of \S\ref{sec:lr-controlled} and the null-handling analysis
of \S\ref{sec:lr-eval}. They differ from FRs and NFRs in being
properties of the \textit{experiment}: a system can be functionally
perfect and still fail ER1, for example by feeding the two
modalities different post-processing, and so invalidate the
comparison.

\begin{figure}[ht]
\centering
\includegraphics[width=\linewidth]{Chapters/figures/fig-4-4-er-table.png}
\caption[Experimental requirements.]{Experimental requirements,
each traced to the literature-review subsection that motivates it.}
\label{fig:er}
\end{figure}

Figure~\ref{fig:er} closes the traceability chain \textit{lit-review
subsection $\to$ ER $\to$ design decision $\to$ test $\to$ evaluation
revisit}. ER1--ER7 are all \textit{Must}; any violation invalidates
the study against the rubric this report has committed to.

% ------------------------------------------------------------
\section{Experimental specification}
\label{sec:spec-experimental}

The experiment is a $2\times 2$ factorial over input modality
(OCR-mediated, vision-mediated) and LLM provider (OpenAI,
Anthropic). The four conditions are labelled C1--C4.

\begin{figure}[ht]
\centering
\includegraphics[width=0.78\linewidth]{Chapters/figures/fig-4-1-factorial-matrix.png}
\caption[$2\times 2$ factorial matrix.]{The four conditions under
which the system is run. Rows: LLM provider. Columns: input
modality. Each cell is a complete pipeline run over the full
multilingual utility-bill dataset.}
\label{fig:factorial}
\end{figure}

Within a column, modality is fixed and the reasoning engine changes
(RQ3). Within a row, the reasoning engine is fixed and modality
changes (RQ1). All four cells are used jointly for RQ2.

Each RQ is pre-specified below: metric, baselines, hypothesis,
controls, and success criterion fixed before any run, so the
evaluation chapter only reports whether each condition met or
missed those criteria.

\paragraph{RQ1 pre-specification.} \textit{Metric:} overall,
document-level, and per-field accuracy
(\S\ref{sec:spec-evaluation}), independently for C1--C4.
\textit{Baselines:} a null-prediction baseline returning
\texttt{null} for every field, to establish the floor implied by
the both-null rule; each condition paired against its within-provider
counterpart (C1 vs.~C2, C3 vs.~C4); Shen et al.'s directional prior
that image-only MLLM input can reach parity with OCR-enhanced
setups~\cite{shen2026ocrornot}. \textit{Hypothesis:} H1.
\textit{Controls:} ER1, plus shared schema, normalisation, and
evaluation code. \textit{Success:} the within-provider modality gap
is measurable at field-level granularity and its sign is
attributable via ER5 to a named mechanism (OCR propagation, vision
hallucination, output-format reliability) rather than asserted.

\paragraph{RQ2 pre-specification.} \textit{Metric:} mean per-document
latency with the per-stage breakdown of NFR7, mean per-document
token count, and mean per-document dollar cost, each for C1--C4.
\textit{Baselines:} Tesseract-only latency as the OCR-mediated stage
ceiling; first-two-page vision-mediated cost as a controlled upper
bound on image-token input;
FastVLM~\cite{vasu2024fastvlm} throughput figures as external
reference. \textit{Hypothesis:} H2. \textit{Controls:} ER1 and
identical document set; identical run ordering where API rate limits
permit. \textit{Success:} cost, latency, and operational complexity
are reported as three separable quantities, with no single-winner
collapse without direct quantitative support.

\paragraph{RQ3 pre-specification.} \textit{Metric:} the
within-provider modality gap for each provider, and the difference
between those gaps with its sign. \textit{Baselines:} both providers
on the same dataset, prompt, normalisation, and evaluation harness,
so provider is the only uncontrolled variable across rows.
\textit{Hypothesis:} H3. \textit{Controls:} ER1, ER2, ER6.
\textit{Success:} a provider-specificity claim is accepted only
when the sign or magnitude of the modality gap matches a named
mechanism, for example native constrained
decoding~\cite{openai2024structured} versus prompt-instructed JSON,
or vision-encoder resolution for low-frequency layouts; otherwise
the null is reported.

\paragraph{Slicing.} Accuracy metrics are reported in aggregate and
sliced by \texttt{language} and \texttt{utility\_type}. Slice
accuracies are descriptive given the corpus size; formal
significance testing is deferred (out-of-scope list,
\S\ref{sec:spec-scope}).

% ------------------------------------------------------------
\section{Canonical data schema}
\label{sec:spec-schema}

The extraction target is a 12-field canonical schema
(Figure~\ref{fig:schema}). Fixing the schema is what makes FR5--FR7
objective: an unbounded key-value target has no ground truth against
which hallucination, omission, or value mismatch can be defined, and
the null-aware evaluation (\S\ref{sec:spec-evaluation}) relies on
the closed-world assumption that ``not in the schema'' is not a
valid prediction. All 12 fields are nullable; a given bill may
omit any of them.

\begin{figure}[ht]
\centering
\includegraphics[width=\linewidth]{Chapters/figures/fig-4-5-schema-table.png}
\caption[\texttt{BillExtraction} schema.]{The canonical
\texttt{BillExtraction} schema. All fields are nullable; types are
enforced by Pydantic at parse time.}
\label{fig:schema}
\end{figure}

The 12 fields are the intersection of the minimum-content
requirements mandated for electricity bills by Article~18 of the EU
electricity directive~\cite{eu2019electricity} and the stable core
fields ReceiptSense~\cite{abdallah2024receiptsense} observes across
multilingual receipts. Provider-specific fields outside that core
(per-tariff itemisation, previous-balance reconciliation) are
excluded to keep the comparison constant across issuers whose
layouts differ substantially in those areas.

\paragraph{Normalisation.} Provider-name legal-suffix stripping
(\textit{Ltd}, \textit{GmbH}, \textit{S.p.A.}); EU- and US-style
decimal-separator parsing; ISO~8601 canonicalisation of dates from
locale-specific surface forms; multilingual currency and
consumption-unit synonym mapping (Italian ``SMC'' maps to m$^3$).
Full rules are in Appendix~B.

\paragraph{Ground-truth format.} Ground truth is stored per document
as JSON:

\begin{verbatim}
{
  "document_id": "<unique key>",
  "annotated_by": "<annotator>",
  "annotation_date": "<ISO 8601>",
  "fields": {
    "provider_name": "<value or null>",
    ...
  }
}
\end{verbatim}

Fields absent from the bill are set to \texttt{null} explicitly,
distinguishing ``absent'' from ``unannotated''. The four-class null
taxonomy (\S\ref{sec:spec-evaluation}) depends on this being decided
at annotation time rather than inferred from missing keys.

% ------------------------------------------------------------
\section{Evaluation criteria}
\label{sec:spec-evaluation}

Criteria are fixed here so Chapter~\ref{chap:evaluation} reports
against a target chosen before the data is seen. Executable decision
flows for null classification and failure diagnosis are deferred to
Chapter~\ref{chap:design} \S\ref{sec:design-evalframework}.

\paragraph{Field-level accuracy.} For each field~$f$, let $C_f$,
$I_f$, $H_f$, $O_f$, $B_f$ be the counts across all documents of
correct, incorrect, hallucinated, omitted, and both-null predictions.
Field accuracy is
\(A_f = C_f / (C_f + I_f + H_f + O_f)\). The both-null count $B_f$ is
excluded from the denominator because including it would inflate
accuracy by rewarding legitimate absences, an effect flagged by
Chinchor's MUC-7 response-key vocabulary~\cite{chinchor1998muc7}
and by ExtractBench~\cite{ferguson2026extractbench}.

\paragraph{Document-level accuracy.} The fraction of documents with
zero incorrect, hallucinated, or omitted fields across all 12
targets. A single wrong field invalidates the extraction for
downstream bookkeeping, so this is the most production-relevant
measure.

\paragraph{Overall accuracy.} Field-level counts aggregated across
all 12 fields and all documents under the same both-null-excluded
denominator.

\paragraph{Four-class null taxonomy.} Every field comparison falls
into one of the five cells of Figure~\ref{fig:null-taxonomy}. The
distinction between hallucination and omission follows Dale et
al.~\cite{dale2023hallucination} and is required by ER6.

\begin{figure}[ht]
\centering
\includegraphics[width=0.85\linewidth]{Chapters/figures/fig-4-6-null-taxonomy.png}
\caption[Null-aware taxonomy.]{The null-aware evaluation taxonomy
applied to every field comparison.}
\label{fig:null-taxonomy}
\end{figure}

\paragraph{Comparison rules.} Floats (\texttt{total\_amount\_due},
\texttt{consumption\_amount}) compare with absolute tolerance
$\pm 0.01$ after normalisation. Dates compare by exact string
equality on the ISO~8601 canonical form. Strings compare on the
normalised surface form (NFC Unicode, lowercase,
whitespace-collapsed); Levenshtein similarity is computed for
diagnosis but does not contribute to accuracy.

\paragraph{Failure attribution, OCR-mediated (ER5).} For each
non-correct field: if the ground-truth value occurs as a
case-insensitive substring of the stored OCR text, the failure is
attributed to the LLM (the signal was available and it misread it);
otherwise the failure is attributed to Tesseract (the signal was
lost before the LLM saw it). The heuristic has known false-positive
and false-negative modes, especially for dates where the same digits
recur in unrelated contexts; its limitations are discussed in
Chapter~\ref{chap:design}. The heuristic supplies an automated,
reproducible, first-order attribution that the OCR-cascade
literature (\S\ref{sec:lr-ocrcascade}) does not.

\paragraph{Failure attribution, vision-mediated.} All non-correct
fields are attributed to the vision-mediated model without finer
granularity, because no intermediate artefact analogous to OCR text
exists. The asymmetry between the two conditions is itself a finding
of the study; Chapter~\ref{chap:evaluation} discusses it rather than
papering over it with a weaker heuristic.

\paragraph{Chart set (FR8).} Six canonical figures are produced by
the reporting stage: (1) overall-accuracy bar chart across
conditions; (2) per-field-accuracy heatmap; (3) per-stage timing
stacked bar chart; (4) cost-per-document bar chart;
(5) accuracy-by-language grouped bar chart; (6)
accuracy-by-\texttt{utility\_type} grouped bar chart. Each supports
one or more of RQ1--RQ3; the chart-to-RQ mapping is given in
Chapter~\ref{chap:design}.

% ------------------------------------------------------------
\section{Summary and forward bridge}
\label{sec:spec-bridge}

This chapter has fixed RQ1--RQ3 and their directional hypotheses,
scope and terminology, FR1--FR10, NFR1--NFR8, and ER1--ER7 with
upstream traceability to Chapter~\ref{chap:litreview}, the $2\times 2$
experimental matrix and per-RQ success criteria, the 12-field
schema, and the null-aware evaluation metrics and failure-attribution
protocol.

Chapter~\ref{chap:design} takes this specification as input. It
states the architectural style chosen to satisfy FR1--FR10 under
NFR1--NFR8 and ER1--ER7, names the ten design decisions that
realise those requirements, identifies the alternative considered
for each, and gives the technical reason each alternative was
rejected. The evaluation framework specified declaratively here is
given its decision-tree realisation in
Chapter~\ref{chap:design}~\S\ref{sec:design-evalframework}.
Chapter~\ref{chap:evaluation} closes the chain by revisiting every
requirement against the run evidence.
.
%  Do NOT add \documentclass or \begin{document} here.
%  Follows Chapter 3 (Literature Review); precedes Chapter 5 (Design).
%  Scope: what the project aims to learn (objectives, RQs, hypotheses)
%  and what the system must satisfy (FR, NFR, ER, experimental
%  specification, schema, evaluation). Architecture and alternatives
%  considered belong to Chapter~\ref{chap:design}.
% ============================================================

\chapter{Objectives, requirements and specification}
\label{chap:spec}

This chapter fixes what the project aims to learn and what the
system must satisfy for that learning to count as evidence. Research
objectives and hypotheses are in \S\ref{sec:spec-objectives}, scope
in \S\ref{sec:spec-scope}, requirements in
\S\S\ref{sec:spec-fr}--\ref{sec:spec-experimental-requirements},
experimental specification in \S\ref{sec:spec-experimental}, the
schema in \S\ref{sec:spec-schema}, and evaluation criteria in
\S\ref{sec:spec-evaluation}. Every requirement traces back to a
subsection of Chapter~\ref{chap:litreview}; how each is realised
architecturally is deferred to Chapter~\ref{chap:design}. This
chapter states \textit{what} and \textit{why}; the next states
\textit{how} and \textit{what was rejected}.

% ------------------------------------------------------------
\section{Research objectives and hypotheses}
\label{sec:spec-objectives}

Three research questions follow from the gap in
\S\ref{sec:lr-gap}. Each is paired with a directional hypothesis
fixed before the run, so a null result stays informative against a
pre-registered position.

\paragraph{RQ1 --- Modality effect.} Does explicit OCR preprocessing
improve or degrade structured-extraction accuracy compared to
direct vision-mediated extraction by the same LLM, on multilingual,
out-of-distribution utility bills?

\textit{Hypothesis H1.} Vision-mediated extraction will match or
exceed OCR-mediated extraction in aggregate accuracy, with the gap
widening on non-Latin-script and low-frequency layouts, because the
OCR error-propagation ceiling identified by
Lopresti~\cite{lopresti2009ocrerrors} and Nguyen et
al.~\cite{nguyen2021postocr} binds on multilingual input and is not
shared by the vision path.

\paragraph{RQ2 --- Operational trade-off.} How do the OCR-mediated
and vision-mediated pipelines compare on per-document monetary
cost, wall-clock latency, and operational complexity?

\textit{Hypothesis H2.} Vision-mediated extraction will cost more
per document in API tokens while eliminating Tesseract language-pack
management, producing a cost-complexity trade-off rather than a
single-winner ranking, consistent with the operational-cost
observations in FastVLM~\cite{vasu2024fastvlm} and
DocVLM~\cite{nacson2024docvlm}.

\paragraph{RQ3 --- Cross-provider generality.} Does the modality
effect observed under RQ1 generalise across LLM providers from
different architectural families, or is it provider-specific?

\textit{Hypothesis H3.} The OCR-vs-vision gap will vary by provider
because structured-output support, vision-encoder resolution, and
multilingual OCR capability differ between
OpenAI~\cite{openai2024gpt4o,openai2024structured} and
Anthropic~\cite{anthropic2024claude} models. A provider-independent
gap would falsify H3.

% ------------------------------------------------------------
\section{Scope, terminology and success framing}
\label{sec:spec-scope}

\paragraph{In scope.} Template-free structured extraction from
multilingual, born-digital and scanned utility bills. Two
general-purpose LLM providers held constant across modalities.
Zero-shot prompting with temperature~$0.0$. Full-document accuracy,
cost, and latency under the null-aware evaluation of
\S\ref{sec:spec-evaluation}. Per-stage attribution of OCR-mediated
failures to OCR text or LLM reasoning.

\paragraph{Out of scope.} Specialised document architectures
(Donut~\cite{kim2022donut}, LayoutLMv3~\cite{huang2022layoutlmv3});
fine-tuning; few-shot prompting; image preprocessing beyond
rasterisation; asynchronous or batched submission; confidence
calibration; significance testing on slice accuracy (sample size
does not support it reliably). A specialised architecture, a
fine-tuning step, or a preprocessing transform would turn the
comparison into a cross-architecture or cross-training study rather
than the single-variable modality test the gap demands. These are
acknowledged as limitations in Chapter~\ref{chap:evaluation}.

\paragraph{Canonical terminology.} The literature uses several names
for each pipeline. For the rest of this report:

\begin{itemize}
\item \textbf{OCR-mediated pipeline}: document image through
Tesseract~\cite{smith2007tesseract}, then the text to the LLM in
text mode. Equivalent to ``OCR-cascade'' and ``text-mode''.
\item \textbf{Vision-mediated pipeline}: document image sent
directly to the LLM in vision mode. Equivalent to ``OCR-free'' and
``vision-mode''.
\item \textbf{Reasoning engine}: the LLM held constant across the
two pipelines within a condition.
\item \textbf{Condition}: a unique (\textit{provider},
\textit{modality}) pair; four exist in the $2\times 2$ factorial,
labelled C1--C4.
\end{itemize}

\paragraph{Success framing.} The project is not chasing a single
leaderboard number. For each RQ, success means a metric collected
under pre-specified conditions, a directional hypothesis fixed
before collection, and a mechanism-level attribution for the sign
of the result. A null finding that is attributable to a named
mechanism (OCR propagation, vision hallucination, output-format
asymmetry, provider-specific encoder limitations) is a successful
outcome; an uninterpretable headline number is not.

% ------------------------------------------------------------
\section{Functional requirements}
\label{sec:spec-fr}

Functional requirements state what the system must do. Each is
verifiable against a run artefact, an automated test, or an
evaluation metric. Priority uses a two-level MoSCoW reduction:
\textit{Must} requirements are those whose violation invalidates
the experiment; \textit{Should} requirements improve operational
usability without affecting validity.

\begin{figure}[ht]
\centering
\includegraphics[width=\linewidth]{Chapters/figures/fig-4-2-fr-table.png}
\caption[Functional requirements catalogue.]{Functional
requirements catalogue. Each requirement names both its
verification artefact and its automated test.}
\label{fig:fr}
\end{figure}

Figure~\ref{fig:fr} fixes visible behaviour, not internal structure;
how ingestion, pipelines, and evaluation are organised inside the
system is covered in Chapter~\ref{chap:design}.

% ------------------------------------------------------------
\section{Non-functional requirements}
\label{sec:spec-nfr}

Non-functional requirements state properties the system must exhibit
irrespective of any one extraction. They guard experimental validity
against engineering failure modes documented in the literature:
irreproducibility, output-mechanism asymmetry, silent configuration
drift, and runaway
cost~\cite{pineau2021reproducibility,openai2024structured}. A system
can satisfy an FR while violating an NFR; such a system produces
extractions but not an experiment.

\begin{figure}[ht]
\centering
\includegraphics[width=\linewidth]{Chapters/figures/fig-4-3-nfr-table.png}
\caption[Non-functional requirements catalogue.]{Non-functional
requirements catalogue. All eight are priority \textit{Must}.}
\label{fig:nfr}
\end{figure}

NFR1--NFR8 are all \textit{Must}; each one is a precondition for a
finding to be defensible under cold-reading scrutiny.

% ------------------------------------------------------------
\section{Experimental requirements}
\label{sec:spec-experimental-requirements}

Experimental requirements are the seven classes derived from the
rubric of \S\ref{sec:lr-controlled} and the null-handling analysis
of \S\ref{sec:lr-eval}. They differ from FRs and NFRs in being
properties of the \textit{experiment}: a system can be functionally
perfect and still fail ER1, for example by feeding the two
modalities different post-processing, and so invalidate the
comparison.

\begin{figure}[ht]
\centering
\includegraphics[width=\linewidth]{Chapters/figures/fig-4-4-er-table.png}
\caption[Experimental requirements.]{Experimental requirements,
each traced to the literature-review subsection that motivates it.}
\label{fig:er}
\end{figure}

Figure~\ref{fig:er} closes the traceability chain \textit{lit-review
subsection $\to$ ER $\to$ design decision $\to$ test $\to$ evaluation
revisit}. ER1--ER7 are all \textit{Must}; any violation invalidates
the study against the rubric this report has committed to.

% ------------------------------------------------------------
\section{Experimental specification}
\label{sec:spec-experimental}

The experiment is a $2\times 2$ factorial over input modality
(OCR-mediated, vision-mediated) and LLM provider (OpenAI,
Anthropic). The four conditions are labelled C1--C4.

\begin{figure}[ht]
\centering
\includegraphics[width=0.78\linewidth]{Chapters/figures/fig-4-1-factorial-matrix.png}
\caption[$2\times 2$ factorial matrix.]{The four conditions under
which the system is run. Rows: LLM provider. Columns: input
modality. Each cell is a complete pipeline run over the full
multilingual utility-bill dataset.}
\label{fig:factorial}
\end{figure}

Within a column, modality is fixed and the reasoning engine changes
(RQ3). Within a row, the reasoning engine is fixed and modality
changes (RQ1). All four cells are used jointly for RQ2.

Each RQ is pre-specified below: metric, baselines, hypothesis,
controls, and success criterion fixed before any run, so the
evaluation chapter only reports whether each condition met or
missed those criteria.

\paragraph{RQ1 pre-specification.} \textit{Metric:} overall,
document-level, and per-field accuracy
(\S\ref{sec:spec-evaluation}), independently for C1--C4.
\textit{Baselines:} a null-prediction baseline returning
\texttt{null} for every field, to establish the floor implied by
the both-null rule; each condition paired against its within-provider
counterpart (C1 vs.~C2, C3 vs.~C4); Shen et al.'s directional prior
that image-only MLLM input can reach parity with OCR-enhanced
setups~\cite{shen2026ocrornot}. \textit{Hypothesis:} H1.
\textit{Controls:} ER1, plus shared schema, normalisation, and
evaluation code. \textit{Success:} the within-provider modality gap
is measurable at field-level granularity and its sign is
attributable via ER5 to a named mechanism (OCR propagation, vision
hallucination, output-format reliability) rather than asserted.

\paragraph{RQ2 pre-specification.} \textit{Metric:} mean per-document
latency with the per-stage breakdown of NFR7, mean per-document
token count, and mean per-document dollar cost, each for C1--C4.
\textit{Baselines:} Tesseract-only latency as the OCR-mediated stage
ceiling; first-two-page vision-mediated cost as a controlled upper
bound on image-token input;
FastVLM~\cite{vasu2024fastvlm} throughput figures as external
reference. \textit{Hypothesis:} H2. \textit{Controls:} ER1 and
identical document set; identical run ordering where API rate limits
permit. \textit{Success:} cost, latency, and operational complexity
are reported as three separable quantities, with no single-winner
collapse without direct quantitative support.

\paragraph{RQ3 pre-specification.} \textit{Metric:} the
within-provider modality gap for each provider, and the difference
between those gaps with its sign. \textit{Baselines:} both providers
on the same dataset, prompt, normalisation, and evaluation harness,
so provider is the only uncontrolled variable across rows.
\textit{Hypothesis:} H3. \textit{Controls:} ER1, ER2, ER6.
\textit{Success:} a provider-specificity claim is accepted only
when the sign or magnitude of the modality gap matches a named
mechanism, for example native constrained
decoding~\cite{openai2024structured} versus prompt-instructed JSON,
or vision-encoder resolution for low-frequency layouts; otherwise
the null is reported.

\paragraph{Slicing.} Accuracy metrics are reported in aggregate and
sliced by \texttt{language} and \texttt{utility\_type}. Slice
accuracies are descriptive given the corpus size; formal
significance testing is deferred (out-of-scope list,
\S\ref{sec:spec-scope}).

% ------------------------------------------------------------
\section{Canonical data schema}
\label{sec:spec-schema}

The extraction target is a 12-field canonical schema
(Figure~\ref{fig:schema}). Fixing the schema is what makes FR5--FR7
objective: an unbounded key-value target has no ground truth against
which hallucination, omission, or value mismatch can be defined, and
the null-aware evaluation (\S\ref{sec:spec-evaluation}) relies on
the closed-world assumption that ``not in the schema'' is not a
valid prediction. All 12 fields are nullable; a given bill may
omit any of them.

\begin{figure}[ht]
\centering
\includegraphics[width=\linewidth]{Chapters/figures/fig-4-5-schema-table.png}
\caption[\texttt{BillExtraction} schema.]{The canonical
\texttt{BillExtraction} schema. All fields are nullable; types are
enforced by Pydantic at parse time.}
\label{fig:schema}
\end{figure}

The 12 fields are the intersection of the minimum-content
requirements mandated for electricity bills by Article~18 of the EU
electricity directive~\cite{eu2019electricity} and the stable core
fields ReceiptSense~\cite{abdallah2024receiptsense} observes across
multilingual receipts. Provider-specific fields outside that core
(per-tariff itemisation, previous-balance reconciliation) are
excluded to keep the comparison constant across issuers whose
layouts differ substantially in those areas.

\paragraph{Normalisation.} Provider-name legal-suffix stripping
(\textit{Ltd}, \textit{GmbH}, \textit{S.p.A.}); EU- and US-style
decimal-separator parsing; ISO~8601 canonicalisation of dates from
locale-specific surface forms; multilingual currency and
consumption-unit synonym mapping (Italian ``SMC'' maps to m$^3$).
Full rules are in Appendix~B.

\paragraph{Ground-truth format.} Ground truth is stored per document
as JSON:

\begin{verbatim}
{
  "document_id": "<unique key>",
  "annotated_by": "<annotator>",
  "annotation_date": "<ISO 8601>",
  "fields": {
    "provider_name": "<value or null>",
    ...
  }
}
\end{verbatim}

Fields absent from the bill are set to \texttt{null} explicitly,
distinguishing ``absent'' from ``unannotated''. The four-class null
taxonomy (\S\ref{sec:spec-evaluation}) depends on this being decided
at annotation time rather than inferred from missing keys.

% ------------------------------------------------------------
\section{Evaluation criteria}
\label{sec:spec-evaluation}

Criteria are fixed here so Chapter~\ref{chap:evaluation} reports
against a target chosen before the data is seen. Executable decision
flows for null classification and failure diagnosis are deferred to
Chapter~\ref{chap:design} \S\ref{sec:design-evalframework}.

\paragraph{Field-level accuracy.} For each field~$f$, let $C_f$,
$I_f$, $H_f$, $O_f$, $B_f$ be the counts across all documents of
correct, incorrect, hallucinated, omitted, and both-null predictions.
Field accuracy is
\(A_f = C_f / (C_f + I_f + H_f + O_f)\). The both-null count $B_f$ is
excluded from the denominator because including it would inflate
accuracy by rewarding legitimate absences, an effect flagged by
Chinchor's MUC-7 response-key vocabulary~\cite{chinchor1998muc7}
and by ExtractBench~\cite{ferguson2026extractbench}.

\paragraph{Document-level accuracy.} The fraction of documents with
zero incorrect, hallucinated, or omitted fields across all 12
targets. A single wrong field invalidates the extraction for
downstream bookkeeping, so this is the most production-relevant
measure.

\paragraph{Overall accuracy.} Field-level counts aggregated across
all 12 fields and all documents under the same both-null-excluded
denominator.

\paragraph{Four-class null taxonomy.} Every field comparison falls
into one of the five cells of Figure~\ref{fig:null-taxonomy}. The
distinction between hallucination and omission follows Dale et
al.~\cite{dale2023hallucination} and is required by ER6.

\begin{figure}[ht]
\centering
\includegraphics[width=0.85\linewidth]{Chapters/figures/fig-4-6-null-taxonomy.png}
\caption[Null-aware taxonomy.]{The null-aware evaluation taxonomy
applied to every field comparison.}
\label{fig:null-taxonomy}
\end{figure}

\paragraph{Comparison rules.} Floats (\texttt{total\_amount\_due},
\texttt{consumption\_amount}) compare with absolute tolerance
$\pm 0.01$ after normalisation. Dates compare by exact string
equality on the ISO~8601 canonical form. Strings compare on the
normalised surface form (NFC Unicode, lowercase,
whitespace-collapsed); Levenshtein similarity is computed for
diagnosis but does not contribute to accuracy.

\paragraph{Failure attribution, OCR-mediated (ER5).} For each
non-correct field: if the ground-truth value occurs as a
case-insensitive substring of the stored OCR text, the failure is
attributed to the LLM (the signal was available and it misread it);
otherwise the failure is attributed to Tesseract (the signal was
lost before the LLM saw it). The heuristic has known false-positive
and false-negative modes, especially for dates where the same digits
recur in unrelated contexts; its limitations are discussed in
Chapter~\ref{chap:design}. The heuristic supplies an automated,
reproducible, first-order attribution that the OCR-cascade
literature (\S\ref{sec:lr-ocrcascade}) does not.

\paragraph{Failure attribution, vision-mediated.} All non-correct
fields are attributed to the vision-mediated model without finer
granularity, because no intermediate artefact analogous to OCR text
exists. The asymmetry between the two conditions is itself a finding
of the study; Chapter~\ref{chap:evaluation} discusses it rather than
papering over it with a weaker heuristic.

\paragraph{Chart set (FR8).} Six canonical figures are produced by
the reporting stage: (1) overall-accuracy bar chart across
conditions; (2) per-field-accuracy heatmap; (3) per-stage timing
stacked bar chart; (4) cost-per-document bar chart;
(5) accuracy-by-language grouped bar chart; (6)
accuracy-by-\texttt{utility\_type} grouped bar chart. Each supports
one or more of RQ1--RQ3; the chart-to-RQ mapping is given in
Chapter~\ref{chap:design}.

% ------------------------------------------------------------
\section{Summary and forward bridge}
\label{sec:spec-bridge}

This chapter has fixed RQ1--RQ3 and their directional hypotheses,
scope and terminology, FR1--FR10, NFR1--NFR8, and ER1--ER7 with
upstream traceability to Chapter~\ref{chap:litreview}, the $2\times 2$
experimental matrix and per-RQ success criteria, the 12-field
schema, and the null-aware evaluation metrics and failure-attribution
protocol.

Chapter~\ref{chap:design} takes this specification as input. It
states the architectural style chosen to satisfy FR1--FR10 under
NFR1--NFR8 and ER1--ER7, names the ten design decisions that
realise those requirements, identifies the alternative considered
for each, and gives the technical reason each alternative was
rejected. The evaluation framework specified declaratively here is
given its decision-tree realisation in
Chapter~\ref{chap:design}~\S\ref{sec:design-evalframework}.
Chapter~\ref{chap:evaluation} closes the chain by revisiting every
requirement against the run evidence.
.
%  Do NOT add \documentclass or \begin{document} here.
%  Follows Chapter 3 (Literature Review); precedes Chapter 5 (Design).
%  Scope: what the project aims to learn (objectives, RQs, hypotheses)
%  and what the system must satisfy (FR, NFR, ER, experimental
%  specification, schema, evaluation). Architecture and alternatives
%  considered belong to Chapter~\ref{chap:design}.
% ============================================================

\chapter{Objectives, requirements and specification}
\label{chap:spec}

This chapter fixes what the project aims to learn and what the
system must satisfy for that learning to count as evidence. Research
objectives and hypotheses are in \S\ref{sec:spec-objectives}, scope
in \S\ref{sec:spec-scope}, requirements in
\S\S\ref{sec:spec-fr}--\ref{sec:spec-experimental-requirements},
experimental specification in \S\ref{sec:spec-experimental}, the
schema in \S\ref{sec:spec-schema}, and evaluation criteria in
\S\ref{sec:spec-evaluation}. Every requirement traces back to a
subsection of Chapter~\ref{chap:litreview}; how each is realised
architecturally is deferred to Chapter~\ref{chap:design}. This
chapter states \textit{what} and \textit{why}; the next states
\textit{how} and \textit{what was rejected}.

% ------------------------------------------------------------
\section{Research objectives and hypotheses}
\label{sec:spec-objectives}

Three research questions follow from the gap in
\S\ref{sec:lr-gap}. Each is paired with a directional hypothesis
fixed before the run, so a null result stays informative against a
pre-registered position.

\paragraph{RQ1 --- Modality effect.} Does explicit OCR preprocessing
improve or degrade structured-extraction accuracy compared to
direct vision-mediated extraction by the same LLM, on multilingual,
out-of-distribution utility bills?

\textit{Hypothesis H1.} Vision-mediated extraction will match or
exceed OCR-mediated extraction in aggregate accuracy, with the gap
widening on non-Latin-script and low-frequency layouts, because the
OCR error-propagation ceiling identified by
Lopresti~\cite{lopresti2009ocrerrors} and Nguyen et
al.~\cite{nguyen2021postocr} binds on multilingual input and is not
shared by the vision path.

\paragraph{RQ2 --- Operational trade-off.} How do the OCR-mediated
and vision-mediated pipelines compare on per-document monetary
cost, wall-clock latency, and operational complexity?

\textit{Hypothesis H2.} Vision-mediated extraction will cost more
per document in API tokens while eliminating Tesseract language-pack
management, producing a cost-complexity trade-off rather than a
single-winner ranking, consistent with the operational-cost
observations in FastVLM~\cite{vasu2024fastvlm} and
DocVLM~\cite{nacson2024docvlm}.

\paragraph{RQ3 --- Cross-provider generality.} Does the modality
effect observed under RQ1 generalise across LLM providers from
different architectural families, or is it provider-specific?

\textit{Hypothesis H3.} The OCR-vs-vision gap will vary by provider
because structured-output support, vision-encoder resolution, and
multilingual OCR capability differ between
OpenAI~\cite{openai2024gpt4o,openai2024structured} and
Anthropic~\cite{anthropic2024claude} models. A provider-independent
gap would falsify H3.

% ------------------------------------------------------------
\section{Scope, terminology and success framing}
\label{sec:spec-scope}

\paragraph{In scope.} Template-free structured extraction from
multilingual, born-digital and scanned utility bills. Two
general-purpose LLM providers held constant across modalities.
Zero-shot prompting with temperature~$0.0$. Full-document accuracy,
cost, and latency under the null-aware evaluation of
\S\ref{sec:spec-evaluation}. Per-stage attribution of OCR-mediated
failures to OCR text or LLM reasoning.

\paragraph{Out of scope.} Specialised document architectures
(Donut~\cite{kim2022donut}, LayoutLMv3~\cite{huang2022layoutlmv3});
fine-tuning; few-shot prompting; image preprocessing beyond
rasterisation; asynchronous or batched submission; confidence
calibration; significance testing on slice accuracy (sample size
does not support it reliably). A specialised architecture, a
fine-tuning step, or a preprocessing transform would turn the
comparison into a cross-architecture or cross-training study rather
than the single-variable modality test the gap demands. These are
acknowledged as limitations in Chapter~\ref{chap:evaluation}.

\paragraph{Canonical terminology.} The literature uses several names
for each pipeline. For the rest of this report:

\begin{itemize}
\item \textbf{OCR-mediated pipeline}: document image through
Tesseract~\cite{smith2007tesseract}, then the text to the LLM in
text mode. Equivalent to ``OCR-cascade'' and ``text-mode''.
\item \textbf{Vision-mediated pipeline}: document image sent
directly to the LLM in vision mode. Equivalent to ``OCR-free'' and
``vision-mode''.
\item \textbf{Reasoning engine}: the LLM held constant across the
two pipelines within a condition.
\item \textbf{Condition}: a unique (\textit{provider},
\textit{modality}) pair; four exist in the $2\times 2$ factorial,
labelled C1--C4.
\end{itemize}

\paragraph{Success framing.} The project is not chasing a single
leaderboard number. For each RQ, success means a metric collected
under pre-specified conditions, a directional hypothesis fixed
before collection, and a mechanism-level attribution for the sign
of the result. A null finding that is attributable to a named
mechanism (OCR propagation, vision hallucination, output-format
asymmetry, provider-specific encoder limitations) is a successful
outcome; an uninterpretable headline number is not.

% ------------------------------------------------------------
\section{Functional requirements}
\label{sec:spec-fr}

Functional requirements state what the system must do. Each is
verifiable against a run artefact, an automated test, or an
evaluation metric. Priority uses a two-level MoSCoW reduction:
\textit{Must} requirements are those whose violation invalidates
the experiment; \textit{Should} requirements improve operational
usability without affecting validity.

\begin{figure}[ht]
\centering
\includegraphics[width=\linewidth]{Chapters/figures/fig-4-2-fr-table.png}
\caption[Functional requirements catalogue.]{Functional
requirements catalogue. Each requirement names both its
verification artefact and its automated test.}
\label{fig:fr}
\end{figure}

Figure~\ref{fig:fr} fixes visible behaviour, not internal structure;
how ingestion, pipelines, and evaluation are organised inside the
system is covered in Chapter~\ref{chap:design}.

% ------------------------------------------------------------
\section{Non-functional requirements}
\label{sec:spec-nfr}

Non-functional requirements state properties the system must exhibit
irrespective of any one extraction. They guard experimental validity
against engineering failure modes documented in the literature:
irreproducibility, output-mechanism asymmetry, silent configuration
drift, and runaway
cost~\cite{pineau2021reproducibility,openai2024structured}. A system
can satisfy an FR while violating an NFR; such a system produces
extractions but not an experiment.

\begin{figure}[ht]
\centering
\includegraphics[width=\linewidth]{Chapters/figures/fig-4-3-nfr-table.png}
\caption[Non-functional requirements catalogue.]{Non-functional
requirements catalogue. All eight are priority \textit{Must}.}
\label{fig:nfr}
\end{figure}

NFR1--NFR8 are all \textit{Must}; each one is a precondition for a
finding to be defensible under cold-reading scrutiny.

% ------------------------------------------------------------
\section{Experimental requirements}
\label{sec:spec-experimental-requirements}

Experimental requirements are the seven classes derived from the
rubric of \S\ref{sec:lr-controlled} and the null-handling analysis
of \S\ref{sec:lr-eval}. They differ from FRs and NFRs in being
properties of the \textit{experiment}: a system can be functionally
perfect and still fail ER1, for example by feeding the two
modalities different post-processing, and so invalidate the
comparison.

\begin{figure}[ht]
\centering
\includegraphics[width=\linewidth]{Chapters/figures/fig-4-4-er-table.png}
\caption[Experimental requirements.]{Experimental requirements,
each traced to the literature-review subsection that motivates it.}
\label{fig:er}
\end{figure}

Figure~\ref{fig:er} closes the traceability chain \textit{lit-review
subsection $\to$ ER $\to$ design decision $\to$ test $\to$ evaluation
revisit}. ER1--ER7 are all \textit{Must}; any violation invalidates
the study against the rubric this report has committed to.

% ------------------------------------------------------------
\section{Experimental specification}
\label{sec:spec-experimental}

The experiment is a $2\times 2$ factorial over input modality
(OCR-mediated, vision-mediated) and LLM provider (OpenAI,
Anthropic). The four conditions are labelled C1--C4.

\begin{figure}[ht]
\centering
\includegraphics[width=0.78\linewidth]{Chapters/figures/fig-4-1-factorial-matrix.png}
\caption[$2\times 2$ factorial matrix.]{The four conditions under
which the system is run. Rows: LLM provider. Columns: input
modality. Each cell is a complete pipeline run over the full
multilingual utility-bill dataset.}
\label{fig:factorial}
\end{figure}

Within a column, modality is fixed and the reasoning engine changes
(RQ3). Within a row, the reasoning engine is fixed and modality
changes (RQ1). All four cells are used jointly for RQ2.

Each RQ is pre-specified below: metric, baselines, hypothesis,
controls, and success criterion fixed before any run, so the
evaluation chapter only reports whether each condition met or
missed those criteria.

\paragraph{RQ1 pre-specification.} \textit{Metric:} overall,
document-level, and per-field accuracy
(\S\ref{sec:spec-evaluation}), independently for C1--C4.
\textit{Baselines:} a null-prediction baseline returning
\texttt{null} for every field, to establish the floor implied by
the both-null rule; each condition paired against its within-provider
counterpart (C1 vs.~C2, C3 vs.~C4); Shen et al.'s directional prior
that image-only MLLM input can reach parity with OCR-enhanced
setups~\cite{shen2026ocrornot}. \textit{Hypothesis:} H1.
\textit{Controls:} ER1, plus shared schema, normalisation, and
evaluation code. \textit{Success:} the within-provider modality gap
is measurable at field-level granularity and its sign is
attributable via ER5 to a named mechanism (OCR propagation, vision
hallucination, output-format reliability) rather than asserted.

\paragraph{RQ2 pre-specification.} \textit{Metric:} mean per-document
latency with the per-stage breakdown of NFR7, mean per-document
token count, and mean per-document dollar cost, each for C1--C4.
\textit{Baselines:} Tesseract-only latency as the OCR-mediated stage
ceiling; first-two-page vision-mediated cost as a controlled upper
bound on image-token input;
FastVLM~\cite{vasu2024fastvlm} throughput figures as external
reference. \textit{Hypothesis:} H2. \textit{Controls:} ER1 and
identical document set; identical run ordering where API rate limits
permit. \textit{Success:} cost, latency, and operational complexity
are reported as three separable quantities, with no single-winner
collapse without direct quantitative support.

\paragraph{RQ3 pre-specification.} \textit{Metric:} the
within-provider modality gap for each provider, and the difference
between those gaps with its sign. \textit{Baselines:} both providers
on the same dataset, prompt, normalisation, and evaluation harness,
so provider is the only uncontrolled variable across rows.
\textit{Hypothesis:} H3. \textit{Controls:} ER1, ER2, ER6.
\textit{Success:} a provider-specificity claim is accepted only
when the sign or magnitude of the modality gap matches a named
mechanism, for example native constrained
decoding~\cite{openai2024structured} versus prompt-instructed JSON,
or vision-encoder resolution for low-frequency layouts; otherwise
the null is reported.

\paragraph{Slicing.} Accuracy metrics are reported in aggregate and
sliced by \texttt{language} and \texttt{utility\_type}. Slice
accuracies are descriptive given the corpus size; formal
significance testing is deferred (out-of-scope list,
\S\ref{sec:spec-scope}).

% ------------------------------------------------------------
\section{Canonical data schema}
\label{sec:spec-schema}

The extraction target is a 12-field canonical schema
(Figure~\ref{fig:schema}). Fixing the schema is what makes FR5--FR7
objective: an unbounded key-value target has no ground truth against
which hallucination, omission, or value mismatch can be defined, and
the null-aware evaluation (\S\ref{sec:spec-evaluation}) relies on
the closed-world assumption that ``not in the schema'' is not a
valid prediction. All 12 fields are nullable; a given bill may
omit any of them.

\begin{figure}[ht]
\centering
\includegraphics[width=\linewidth]{Chapters/figures/fig-4-5-schema-table.png}
\caption[\texttt{BillExtraction} schema.]{The canonical
\texttt{BillExtraction} schema. All fields are nullable; types are
enforced by Pydantic at parse time.}
\label{fig:schema}
\end{figure}

The 12 fields are the intersection of the minimum-content
requirements mandated for electricity bills by Article~18 of the EU
electricity directive~\cite{eu2019electricity} and the stable core
fields ReceiptSense~\cite{abdallah2024receiptsense} observes across
multilingual receipts. Provider-specific fields outside that core
(per-tariff itemisation, previous-balance reconciliation) are
excluded to keep the comparison constant across issuers whose
layouts differ substantially in those areas.

\paragraph{Normalisation.} Provider-name legal-suffix stripping
(\textit{Ltd}, \textit{GmbH}, \textit{S.p.A.}); EU- and US-style
decimal-separator parsing; ISO~8601 canonicalisation of dates from
locale-specific surface forms; multilingual currency and
consumption-unit synonym mapping (Italian ``SMC'' maps to m$^3$).
Full rules are in Appendix~B.

\paragraph{Ground-truth format.} Ground truth is stored per document
as JSON:

\begin{verbatim}
{
  "document_id": "<unique key>",
  "annotated_by": "<annotator>",
  "annotation_date": "<ISO 8601>",
  "fields": {
    "provider_name": "<value or null>",
    ...
  }
}
\end{verbatim}

Fields absent from the bill are set to \texttt{null} explicitly,
distinguishing ``absent'' from ``unannotated''. The four-class null
taxonomy (\S\ref{sec:spec-evaluation}) depends on this being decided
at annotation time rather than inferred from missing keys.

% ------------------------------------------------------------
\section{Evaluation criteria}
\label{sec:spec-evaluation}

Criteria are fixed here so Chapter~\ref{chap:evaluation} reports
against a target chosen before the data is seen. Executable decision
flows for null classification and failure diagnosis are deferred to
Chapter~\ref{chap:design} \S\ref{sec:design-evalframework}.

\paragraph{Field-level accuracy.} For each field~$f$, let $C_f$,
$I_f$, $H_f$, $O_f$, $B_f$ be the counts across all documents of
correct, incorrect, hallucinated, omitted, and both-null predictions.
Field accuracy is
\(A_f = C_f / (C_f + I_f + H_f + O_f)\). The both-null count $B_f$ is
excluded from the denominator because including it would inflate
accuracy by rewarding legitimate absences, an effect flagged by
Chinchor's MUC-7 response-key vocabulary~\cite{chinchor1998muc7}
and by ExtractBench~\cite{ferguson2026extractbench}.

\paragraph{Document-level accuracy.} The fraction of documents with
zero incorrect, hallucinated, or omitted fields across all 12
targets. A single wrong field invalidates the extraction for
downstream bookkeeping, so this is the most production-relevant
measure.

\paragraph{Overall accuracy.} Field-level counts aggregated across
all 12 fields and all documents under the same both-null-excluded
denominator.

\paragraph{Four-class null taxonomy.} Every field comparison falls
into one of the five cells of Figure~\ref{fig:null-taxonomy}. The
distinction between hallucination and omission follows Dale et
al.~\cite{dale2023hallucination} and is required by ER6.

\begin{figure}[ht]
\centering
\includegraphics[width=0.85\linewidth]{Chapters/figures/fig-4-6-null-taxonomy.png}
\caption[Null-aware taxonomy.]{The null-aware evaluation taxonomy
applied to every field comparison.}
\label{fig:null-taxonomy}
\end{figure}

\paragraph{Comparison rules.} Floats (\texttt{total\_amount\_due},
\texttt{consumption\_amount}) compare with absolute tolerance
$\pm 0.01$ after normalisation. Dates compare by exact string
equality on the ISO~8601 canonical form. Strings compare on the
normalised surface form (NFC Unicode, lowercase,
whitespace-collapsed); Levenshtein similarity is computed for
diagnosis but does not contribute to accuracy.

\paragraph{Failure attribution, OCR-mediated (ER5).} For each
non-correct field: if the ground-truth value occurs as a
case-insensitive substring of the stored OCR text, the failure is
attributed to the LLM (the signal was available and it misread it);
otherwise the failure is attributed to Tesseract (the signal was
lost before the LLM saw it). The heuristic has known false-positive
and false-negative modes, especially for dates where the same digits
recur in unrelated contexts; its limitations are discussed in
Chapter~\ref{chap:design}. The heuristic supplies an automated,
reproducible, first-order attribution that the OCR-cascade
literature (\S\ref{sec:lr-ocrcascade}) does not.

\paragraph{Failure attribution, vision-mediated.} All non-correct
fields are attributed to the vision-mediated model without finer
granularity, because no intermediate artefact analogous to OCR text
exists. The asymmetry between the two conditions is itself a finding
of the study; Chapter~\ref{chap:evaluation} discusses it rather than
papering over it with a weaker heuristic.

\paragraph{Chart set (FR8).} Six canonical figures are produced by
the reporting stage: (1) overall-accuracy bar chart across
conditions; (2) per-field-accuracy heatmap; (3) per-stage timing
stacked bar chart; (4) cost-per-document bar chart;
(5) accuracy-by-language grouped bar chart; (6)
accuracy-by-\texttt{utility\_type} grouped bar chart. Each supports
one or more of RQ1--RQ3; the chart-to-RQ mapping is given in
Chapter~\ref{chap:design}.

% ------------------------------------------------------------
\section{Summary and forward bridge}
\label{sec:spec-bridge}

This chapter has fixed RQ1--RQ3 and their directional hypotheses,
scope and terminology, FR1--FR10, NFR1--NFR8, and ER1--ER7 with
upstream traceability to Chapter~\ref{chap:litreview}, the $2\times 2$
experimental matrix and per-RQ success criteria, the 12-field
schema, and the null-aware evaluation metrics and failure-attribution
protocol.

Chapter~\ref{chap:design} takes this specification as input. It
states the architectural style chosen to satisfy FR1--FR10 under
NFR1--NFR8 and ER1--ER7, names the ten design decisions that
realise those requirements, identifies the alternative considered
for each, and gives the technical reason each alternative was
rejected. The evaluation framework specified declaratively here is
given its decision-tree realisation in
Chapter~\ref{chap:design}~\S\ref{sec:design-evalframework}.
Chapter~\ref{chap:evaluation} closes the chain by revisiting every
requirement against the run evidence.
